\chapter{几何Satake等价III：最后的证明}\label{L23}\lecturedate{2020/05/09}\sourcepage{191}
仍设$k=\C$，以下遵循\bref{L23MV}{MV}第6--7节。我们处理的情形相对简单。对$k$为一般交换环的更普遍结果，论证较复杂，可参看\bref{L23MV}{MV}第8--12节。
\section{张量函子}\label{L23S1}\idx{张量函子}\idx{融合积}
如\lref{L21}末所述，有自然的函子同构
\begin{equation}\tag{1.1}\label{L23eq11}H^*\cong F=\bigoplus_{\gamma\in X^\vee}F_\gamma:P_{G(\mathcal O)}(\operatorname{Gr},k)\to\{\text{有限生成 }k\text{ 模}\},\end{equation}
其中$F_\gamma=H_c^{2\langle\kappa,\gamma\rangle}(S_\gamma,-)$。在相应的共轭识别下，$F_\gamma$及这一同构不依赖于$T\subseteq B$的选择。因此，整体上同调函子$H^*$忠实且正合。
\originalnumber{1}
\begin{note}[L23S1O1]
虽然\eqref{L23eq11}中的直和遍及无限格$X^\vee$，但对每个反常层$\mathcal A$，$F(\mathcal A)$仍有限生成。按定义，存在正整数$s$，使$\mathcal A$支撑于有限维的有界部分$\operatorname{Gr}^s$，见\eqref{L21eq11}。在$G\hookrightarrow\GL_n$的格坐标中，$\mathcal A|_{S_\gamma}\ne0$蕴含$t^st^\gamma,t^st^{-\gamma}\in M_n(\mathcal O)$，故只可能出现有限多个$\gamma$。
\end{note}
% [Tutor Note: 原书 p191 Remark 1.1 作 $\mathcal A|_{S_\gamma}\ne0\Rightarrow t^\gamma t^s,t^\gamma t^{-s}\in G(\mathcal O)$（第二式即 $\gamma_i\ge s$），与第一式连用只给出单侧约束；按 $\mathcal A$ 支撑于 $\operatorname{Gr}^s$ 应有 $|\gamma_i|\le s$。译文写作 $t^st^{-\gamma}\in M_n(\mathcal O)$，即 $\gamma_i\le s$，正确。]
\sourcepage{192}
由于$k$为域，有限生成$k$模就是有限维$k$向量空间。记其范畴为$\operatorname{Vsp}_k$。以$\operatorname{Vsp}_k(X^\vee)$表示$X^\vee$分次有限维向量空间$V=\bigoplus_{\gamma\in X^\vee}V_\gamma$及保持分次的线性映射所组成的范畴。其张量结构定义为
\[V\otimes W=\bigoplus_{\gamma\in X^\vee}(V\otimes W)_\gamma,\qquad(V\otimes W)_\gamma=\bigoplus_{\gamma_1+\gamma_2=\gamma}V_{\gamma_1}\otimes W_{\gamma_2}.\]
\eqref{L23eq11}中的$F$取值于$\operatorname{Vsp}_k(X^\vee)$。我们已经在$P_{G(\mathcal O)}(\operatorname{Gr},k)$中定义卷积积$*$，使其具有张量范畴结构。
\originalnumber{2}
\begin{proposition}\label{L23S1O2}
\eqref{L23eq11}定义的$F:P_{G(\mathcal O)}(\operatorname{Gr},k)\to\operatorname{Vsp}_k(X^\vee)$为张量函子\znote{张量函子是保持张量结构的函子：除对象上的同构 $F(X\otimes Y)\cong F(X)\otimes F(Y)$ 外，还要求这些同构与结合约束、单位约束相容；本讲中的 $F=H^*$ 把卷积积送到分次向量空间的外张量积。易错点：只在对象上给出同构还不够，必须与约束相容，否则不能谈论张量等价或张量生成。[来源：nLab \textquotedbl monoidal functor\textquotedbl，\url{https://ncatlab.org/nlab/show/monoidal+functor}（Last revised 2026-01-26；2026-11-14 访问）。]}，即
\[F_\gamma(\mathcal A_1*\mathcal A_2)\cong\bigoplus_{\gamma_1+\gamma_2=\gamma}F_{\gamma_1}(\mathcal A_1)\otimes F_{\gamma_2}(\mathcal A_2).\]
\end{proposition}
下面仅给出证明梗概，细节见\bref{L23MV}{MV}。该命题等价于
\begin{equation}\tag{1.2}\label{L23eq12}
H_c^{2\langle\kappa,\gamma\rangle}(S_\gamma,\mathcal A_1*\mathcal A_2)\cong\bigoplus_{\gamma_1+\gamma_2=\gamma}H_c^{2\langle\kappa,\gamma_1\rangle}(S_{\gamma_1},\mathcal A_1)\otimes H_c^{2\langle\kappa,\gamma_2\rangle}(S_{\gamma_2},\mathcal A_2).
\end{equation}
另一方面，有
\begin{equation}\tag{1.3}\label{L23eq13}H_c^m(S_\gamma\times S_\nu,\mathcal A_1\boxtimes\mathcal A_2)\cong\bigoplus_{m_1+m_2=m}H_c^{m_1}(S_\gamma,\mathcal A_1)\otimes H_c^{m_2}(S_\nu,\mathcal A_2).\end{equation}
为了证明\eqref{L23eq12}，需要把仿射Grassmannian“整体化”。回忆\lref{L21}中Beauville–Laszlo对函子$\operatorname{Gr}_G$的解释。设$C$为$k$上的曲线，固定光滑$k$有理点$x\in C$，令$C^*=C\setminus\{x\}$。定义
\[\operatorname{Gr}_{G,x}(R)=\left\{(\mathcal P,\alpha)\ \middle|\ \begin{array}{l}\mathcal P\text{ 是 }C_R\text{ 上的 }G\text{ 丛},\\\alpha:\mathcal P|_{C^*_R}\xrightarrow{\sim}G\times C^*_R\end{array}\right\}\big/\text{同构}.\]
由\eref{定理}{L21S1O4}，它也由$\operatorname{Gr}$代表；为突出所选点$x$，记为$\operatorname{Gr}_x$。考虑函子$\operatorname{Gr}_{G,C^n}:\mathcal R\to\mathbf{Sets}$：
\[\operatorname{Gr}_{G,C^n}(R)=\left\{\begin{array}{l}(x_1,\ldots,x_n)\in C^n(R),\quad\mathcal P\text{ 为 }C_R\text{ 上的 }G\text{ 丛},\\\alpha:\mathcal P|_{C_R\setminus\bigcup_j\Gamma_{x_j}}\xrightarrow{\sim}G\times(C_R\setminus\bigcup_j\Gamma_{x_j})\end{array}\right\}\big/\text{同构},\]
其中$\Gamma_{x_j}$为$x_j$在$C_R$中的图像。
\sourcepage{193}
这一函子由$C^n=C\times\cdots\times C$上的ind概形$\operatorname{Gr}_{C^n}$代表；它在$(x_1,\ldots,x_n)$处的纤维为$\prod_{j=1}^a\operatorname{Gr}_{y_j}$，其中$\{y_1,\ldots,y_a\}=\{x_1,\ldots,x_n\}$，各$y_j$互异。特别地，当$n=2$时，
\[(\operatorname{Gr}_{C^2})_{(x_1,x_2)}\cong\begin{cases}\operatorname{Gr}_{x_1}\times\operatorname{Gr}_{x_2},&x_1\ne x_2,\\\operatorname{Gr}_{x_1},&x_1=x_2.\end{cases}\]
记$\operatorname{Gr}_C=\operatorname{Gr}_{C^1}$。同样可把$G(\mathcal O)$及其等变反常层范畴整体化，得到$C$上的群概形$G_C(\mathcal O)$及范畴
\[P_{G_C(\mathcal O)}(\operatorname{Gr}_C,k).\]
对其中的$\mathcal B_1,\mathcal B_2$，也可定义整体卷积$\mathcal B_1*_C\mathcal B_2$；不过，它位于$\operatorname{Gr}_{C^2}$上，而非$\operatorname{Gr}_C$上。

子空间$S_\gamma,T_\gamma$也能推广到整体Grassmannian。对$\gamma\in X^\vee$，相应的ind子概形$S_\gamma(C^n)$在$(x_1,\ldots,x_n)$处的纤维为
\[\bigsqcup_{\gamma_1+\cdots+\gamma_a=\gamma}\prod_{j=1}^aS_{\gamma_j}\ \subseteq\prod_{j=1}^a\operatorname{Gr}_{y_j}=\pi_n^{-1}(x_1,\ldots,x_n),\]
其中$\pi_n:\operatorname{Gr}_{C^n}\to C^n$为投影。记包含映射为$s_\gamma:S_\gamma(C^n)\hookrightarrow\operatorname{Gr}_{C^n}$。定义$L_\gamma(C^n)$，其对应纤维为
\[\bigsqcup_{\gamma_1+\cdots+\gamma_a=\gamma}\{(L_{\gamma_1},\ldots,L_{\gamma_a})\}\subseteq S_\gamma(C^n).\]
于是有收缩映射$p_\gamma:S_\gamma(C^n)\to L_\gamma(C^n)$。

取$C=\A^1$，则$\operatorname{Gr}_C\cong\operatorname{Gr}\times C$，故有投影$\tau:\operatorname{Gr}_C\to\operatorname{Gr}$。定义$\tau^\circ=\tau^*[1]$；对$\mathcal A\in P_{G(\mathcal O)}(\operatorname{Gr},k)$，$\tau^\circ\mathcal A\in P_{G_C(\mathcal O)}(\operatorname{Gr}_C,k)$。记
\[\mathcal B:=(\tau^\circ\mathcal A_1*_C\tau^\circ\mathcal A_2)[-2].\]
% [Tutor Note: 原书 p193 定义 $\tau^\circ:=\tau^{\mathrm{an}*}[1]$ 之后，直接对 $\tau^\circ\mathcal A_1*_C\tau^\circ\mathcal A_2$ 取 $R^k(\pi_2)_*$；但所列纤维公式相当于对平移 $[-2]$ 后的层取值。译文显式引入 $\mathcal B:=(\tau^\circ\mathcal A_1*_C\tau^\circ\mathcal A_2)[-2]$，以与 p193--p194 的公式一致。]
可证明$C^2$上的层$R^m(\pi_2)_*\mathcal B$为常值层，其纤维为
\[\begin{cases}\displaystyle\bigoplus_{m_1+m_2=m}H^{m_1}(\operatorname{Gr},\mathcal A_1)\otimes H^{m_2}(\operatorname{Gr},\mathcal A_2),&x_1\ne x_2,\\H^m(\operatorname{Gr},\mathcal A_1*\mathcal A_2),&x_1=x_2.\end{cases}\]
于是
\[\bigoplus_{m_1+m_2=m}H^{m_1}(\operatorname{Gr},\mathcal A_1)\otimes H^{m_2}(\operatorname{Gr},\mathcal A_2)\cong H^m(\operatorname{Gr},\mathcal A_1*\mathcal A_2).\]
\sourcepage{194}
此外，有典范分解
\[R^m(\pi_2)_*\mathcal B\cong\bigoplus_{2\langle\kappa,\gamma\rangle=m}\mathcal L_\gamma^m(\mathcal A_1,\mathcal A_2),\]
其中
\[\mathcal L_\gamma^m(\mathcal A_1,\mathcal A_2):=\mathcal H^m\bigl(R(\pi_2|_{L_\gamma})_*R(p_\gamma)_!s_\gamma^*\mathcal B\bigr).\]
因此，每个$\mathcal L_\gamma^m$也为常值层，其纤维为
\[\begin{cases}\displaystyle\bigoplus_{\gamma_1+\gamma_2=\gamma}H_c^m(S_{\gamma_1}\times S_{\gamma_2},\mathcal A_1\boxtimes\mathcal A_2),&x_1\ne x_2,\\H_c^m(S_\gamma,\mathcal A_1*\mathcal A_2),&x_1=x_2.\end{cases}\]
故\eqref{L23eq12}由\eqref{L23eq13}得到。
\section{等变反常层范畴的半单性}\label{L23S2}\idx{奇偶层}
设$Y$为复代数簇，具有分层$Y=\bigsqcup_{\lambda\in\Lambda}Y_\lambda$。记包含映射为$j_\lambda:Y_\lambda\hookrightarrow Y$。
\begin{definition*}
若相对于$\Lambda$可构造的有界$k$层复形$\mathcal F\in D^b_\Lambda(Y,k)$满足
\begin{equation}\tag{2.1}\label{L23eq21}\mathcal H^m(j_\lambda^*\mathcal F)=\mathcal H^m(j_\lambda^!\mathcal F)=0\quad\text{对所有奇数 }m\text{ 及所有 }\lambda,\end{equation}
则称其为\textbf{偶复形}。若$\mathcal F[1]$为偶复形，则称$\mathcal F$为\textbf{奇复形}。若$\mathcal F$可写成$\mathcal F_0\oplus\mathcal F_1$，其中$\mathcal F_0$偶、$\mathcal F_1$奇，则称其为\textbf{奇偶复形}。
\end{definition*}
进一步假设各$Y_\lambda$都单连通，且奇数$m$时$H^m(Y_\lambda,k)=0$。有以下定理。
\originalnumber{1}
\begin{theorem}\label{L23S2O1}
设$\mathcal F$为不可分解的奇偶复形，则：
\begin{enumerate}[label=\arabic*.]
\item $\mathcal F$的支撑不可约，因而等于某个$\overline{Y_\lambda}$；
\item 除去一个平移后，$\mathcal F$在$Y_\lambda$上的限制同构于秩一常值层。
\end{enumerate}
此外，任何两个具有相同支撑的不可分解奇偶复形，在适当平移后同构。
\end{theorem}
\sourcepage{195}
证明见\bref{L23JMW}{JMW}第2.2节。若$\mathcal F$为支撑在$\overline{Y_\lambda}$上的不可分解奇偶复形，则存在唯一的平移，使其为Verdier自对偶；按这一平移归一化的复形称为\textbf{奇偶层}\znote{奇偶层是奇偶复形中的不可分解对象：奇偶复形要求在每层上的限制与余限制的上同调只出现在固定奇偶的次数（式 (2.1)）；对支撑在 $\overline{Y_\lambda}$ 上的不可分解奇偶复形，存在唯一平移使其 Verdier 自对偶，归一化后即为奇偶层。本讲用它把每个 $IC_\lambda$ 的奇偶性固定下来（引理 2.2），这是半单性证明的枢纽。[来源：Juteau--Mautner--Williamson, \textit{Parity sheaves}, J. Amer. Math. Soc. 27 (2014) 1169--1212（本讲参考文献）。]}。

胞腔$\operatorname{Gr}_\lambda$为部分旗簇$G\cdot L_\lambda\cong G/P_\lambda$上的$\C^{n_\lambda}$纤维丛，其中
\[n_\lambda=\sum_{\substack{\alpha\in\Phi\\\langle\alpha,\lambda\rangle>0}}(\langle\alpha,\lambda\rangle-1),\qquad\Delta_\lambda=\{\alpha\text{ 为单根}\mid\langle\alpha,\lambda\rangle=0\},\]
$P_\lambda$为与由$\Delta_\lambda$确定的标准抛物子群相反的抛物子群，参看\sref{L21S2}。
% [Tutor Note: 原书 p195 作 $P_\lambda$ the standard parabolic subgroup given by $\Delta_\lambda$；按本讲左乘 $G(\mathcal K)/G(\mathcal O)$ 的约定，$L_\lambda$ 的稳定子是相应标准抛物子群的反面，译文按此改译（与第 21 讲 p178 的处理一致，跨章统一）。]
Bruhat分解给出$G/P_\lambda$的胞腔分解，从而也给出$\operatorname{Gr}_\lambda$的胞腔分解。特别地，各$\operatorname{Gr}_\lambda$单连通，且没有奇数次数上同调。因此，虽然$\operatorname{Gr}$无限维，\eref{定理}{L23S2O1}仍可在各对象的有限维支撑上用于轨道分层。

回忆，分别以$\mathcal I_!(\lambda,k),\mathcal I_*(\lambda,k),\mathcal I_{!*}(\lambda,k)$表示
\[{}^pj_!(k[\dim\operatorname{Gr}_\lambda]),\quad{}^pj_*(k[\dim\operatorname{Gr}_\lambda]),\quad j_{!*}(k[\dim\operatorname{Gr}_\lambda]),\]
其中$j:\operatorname{Gr}_\lambda\to\operatorname{Gr}$为包含映射。有典范识别
\begin{equation}\tag{2.2}\label{L23eq22}
H_c^{2\langle\kappa,\gamma\rangle}(S_\gamma,\mathcal I_!(\lambda,k))\cong H_c^{2\langle\kappa,\gamma\rangle}(S_\gamma,\mathcal I_*(\lambda,k))\cong k[\operatorname{Irr}(\overline{\operatorname{Gr}_\lambda}\cap S_\gamma)].
\end{equation}
右端表示由$\overline{\operatorname{Gr}_\lambda}\cap S_\gamma$的不可约分支生成的自由$k$模。以下引理可由Kazhdan–Lusztig奇偶性定理推出，见\bref{L23dCM}{dCM}的定理4.4.7或\bref{L23L}{L}第11节；注意这里须使用仿射Weyl群的Kazhdan–Lusztig多项式。
\originalnumber{2}
\begin{lemma}\label{L23S2O2}
反常层$IC_\lambda:=\mathcal I_{!*}(\lambda,k)$具有由$d_\lambda=\dim\operatorname{Gr}_\lambda$决定的奇偶性：$IC_\lambda[d_\lambda]$为偶复形。
\end{lemma}
\originalnumber{3}
\begin{theorem}\label{L23S2O3}范畴$P_{G(\mathcal O)}(\operatorname{Gr},k)$半单\znote{半单性指阿贝尔范畴中每个对象都是单对象的直和，等价地每个短正合列分裂。对 $P_{G(\mathcal O)}(\operatorname{Gr},k)$ 而言，半单性意味着任意反常层都同构于若干 $IC_\lambda$ 的直和，从而其上同调与表示论中的完全可约性一致；这也是定理 3.1 中 $\widetilde G$ 约化的前提之一。易错点：这里半单性并不显然，需要奇偶层理论（引理 2.2）与连通分支的奇偶性论证。[来源：nLab \textquotedbl semisimple category\textquotedbl，\url{https://ncatlab.org/nlab/show/semisimple+category}（Last revised 2023-06-18；2026-11-14 访问）。]}。\end{theorem}
\begin{proof}
由于$P_{G(\mathcal O)}(\operatorname{Gr},k)\cong P_{\mathcal S}(\operatorname{Gr},k)$，由\eref{定理}{L13S3O4}，$IC_\lambda$是这些反常层的简单组成因子。同一连通分支中，各轨道的维数具有相同奇偶性；不同分支间没有扩张。因此，由\eref{引理}{L23S2O2}及限制、余限制的长正合列，每个反常层在每个连通分支上都具有相应的奇偶性。

由\eref{定理}{L23S2O1}，不可分解的反常层必须同构于某个$IC_\lambda$的平移；反常归一化要求这个平移为零。于是每个反常层都为简单$IC_\lambda$的直和，定理得证。
\end{proof}
\originalnumber{4}
\begin{corollary}\label{L23S2O4}$\mathcal I_!(\lambda,k)\cong\mathcal I_*(\lambda,k)\cong IC_\lambda$。\end{corollary}
\begin{proof}
由\eref{定理}{L23S2O3}，$\mathcal I_!(\lambda,k)$与$\mathcal I_*(\lambda,k)$都是$IC_\lambda$及其他$IC_\gamma$的直和。由\eqref{L23eq22}及$j_!$、$j_*$的伴随性质，其间的自然映射为同构。因此，由$IC_\lambda$为这一映射之像的定义，$\mathcal I_!(\lambda,k)\cong\mathcal I_*(\lambda,k)\cong IC_\lambda$。
\end{proof}
\sourcepage{196}
\section{\texorpdfstring{$k=\C$}{k=C}时的最后证明}\label{L23S3}\idx{纤维函子}\idx{Tannaka形式理论}
正合且忠实的张量函子称为\textbf{纤维函子}。因此，在按连通分支奇偶性修正卷积交换约束后，\eqref{L23eq11}中的$F$为纤维函子。
\begin{definition*}
设$\mathcal C$为具有结合、交换和单位约束的张量范畴。若每个对象都有张量对偶，则称$\mathcal C$\textbf{刚性}。此时，对$X,Y\in\mathcal C$，有内部Hom对象$\underline{\Hom}(X,Y)\in\mathcal C$，满足
\[\Hom_{\mathcal C}(Z,\underline{\Hom}(X,Y))\cong\Hom_{\mathcal C}(Z\otimes X,Y),\qquad Z\in\mathcal C.\]
\end{definition*}
对刚性$k$线性阿贝尔张量范畴$\mathcal C$，要求$\End(\mathbf1)=k$。由\bref{L23DM}{DM}的命题1.20和定理2.11，有以下结果。
\originalnumber{1}
\begin{theorem}\label{L23S3O1}
给定纤维函子$F:\mathcal C\to\operatorname{Vsp}_k$，存在仿代数群概形$\widetilde G=\operatorname{Aut}^{\otimes}(F)$，使$F$提升为张量范畴等价$\mathcal C\xrightarrow{\sim}\operatorname{Rep}(\widetilde G)$。此外，当$k$具有特征零且$\mathcal C$半单时，$\widetilde G$为仿约化群。
\end{theorem}
从表示范畴及纤维函子重建代数群的这一方法，通常称为\textbf{Tannaka形式理论}\znote{Tannaka 对偶（Tannaka 重建）断言：带纤维函子 $F:\mathcal C\to\operatorname{Vsp}_k$ 的刚性张量范畴可由自同构群 $\widetilde G=\operatorname{Aut}^{\otimes}(F)$ 重建，且 $\mathcal C\simeq\operatorname{Rep}(\widetilde G)$。直觉是「表示范畴加遗忘函子」与「群」互相决定；定理 3.1 正是这一原理在仿射 Grassmannian 上的实现。易错点：纤维函子的正合性与忠实性是重建的关键假设，缺一不可。[来源：nLab \textquotedbl Tannaka duality\textquotedbl，\url{https://ncatlab.org/nlab/show/Tannaka+duality}（Last revised 2025-07-28；2026-11-14 访问）。]}。
\originalnumber{2}
\begin{note}[L23S3O2]
$P_{G(\mathcal O)}(\operatorname{Gr},k)$的单位对象为支撑在$t^0G(\mathcal O)$上的摩天层$IC_0$。若$\iota$为Cartan对合在$\operatorname{Gr}$上诱导的映射，$\mathbb D$为Verdier对偶，则卷积对偶为$\mathcal F^\vee=\iota^*\mathbb D\mathcal F$，内部Hom为
\[\underline{\Hom}(\mathcal F,\mathcal G)=(\iota^*\mathbb D\mathcal F)*\mathcal G.\]
\end{note}
最后得到以下结论。
\originalnumber{3}
\begin{theorem}\label{L23S3O3}\eref{定理}{L23S3O1}重建出的$\widetilde G$为$G$的Langlands对偶群$G^L$。\end{theorem}
记$T^\vee$为$G^L$的极大环面，则其特征标格$X(T^\vee)=X^\vee$。对每个$\mathcal A$，$F(\mathcal A)=H^*(\operatorname{Gr},\mathcal A)$的$T^\vee$权分解为
\[H^*(\operatorname{Gr},\mathcal A)\cong\bigoplus_{\gamma\in X^\vee}H_c^{2\langle\kappa,\gamma\rangle}(S_\gamma,\mathcal A).\]
给定支配余权$\lambda\in X^\vee$，$F(IC_\lambda)$为最高权$\lambda$的不可约表示。由该定理得到以下推论。
\sourcepage{197}
\originalnumber{4}
\begin{corollary}\label{L23S3O4}
设$V_\lambda$为$G^L$的最高权$\lambda$不可约表示，则其$\gamma$权空间$V_\lambda(\gamma)$可典范识别为由$\overline{\operatorname{Gr}_\lambda}\cap S_\gamma$的不可约分支张成的$k$向量空间。特别地，$\dim V_\lambda(\gamma)$等于这些不可约分支的个数。
\end{corollary}
\begin{proof}[\eref{定理}{L23S3O3}的证明]
首先证明$\widetilde G$为代数群。由\bref{L23DM}{DM}命题2.20，只需证明$\operatorname{Rep}(\widetilde G)$具有一个张量生成元，即存在$V_0$，使每个表示都是某个$V_0^{\otimes n}$的直和项。

以$\Lambda_0$表示支配余权幺半群$(X^\vee)^+$的一组有限生成元，并将$0$加入其中。则$F(\bigoplus_{\lambda\in\Lambda_0}IC_\lambda)$为张量生成元，因为对每个$\gamma\in(X^\vee)^+$，若$\gamma=\sum_{\lambda\in\Lambda_0}a_\lambda\lambda$，$a_\lambda\ge0$，则$IC_\gamma$为$(\bigoplus_{\lambda\in\Lambda_0}IC_\lambda)^{*\sum_\lambda a_\lambda}$的直和项。
% [Tutor Note: 原书 p197 取 $\Delta^\vee$ 为单余根集，并断言每个 $\gamma\in(X^\vee)^+$ 可写成 $\sum a_\lambda\lambda$（$\lambda\in\Delta^\vee$）；但单余根只生成余根格，一般不生成支配余权幺半群。译文改用 $(X^\vee)^+$ 的一组有限生成元 $\Lambda_0$ 并加入 $0$。]

同样，任何含非平凡简单对象$IC_\lambda$的张量子范畴，都必须包含其卷积幂中的$IC_{n\lambda}$，因而具有无限多个不同的简单对象。故不存在由有限多个固定简单对象的直和组成的非平凡张量子范畴。由\bref{L23DM}{DM}推论2.22，$\widetilde G$连通。结合半单性，$\widetilde G$连通且约化。

每个$X^\vee$分次向量空间$V=\bigoplus_\gamma V_\gamma$都可看作特征标格为$X^\vee$的环面$T^\vee$的表示，其中$V_\gamma$为特征标$\gamma:T^\vee\to k^\times$的权空间。带分次的张量函子$\operatorname{Rep}(\widetilde G)\to\operatorname{Vsp}_k(X^\vee)$诱导群同态$h:T^\vee\to\widetilde G$。每个支配权$\lambda$都出现在$F_\lambda(IC_\lambda)$中，而所有支配权生成整个格$X^\vee$，故$h$的核平凡。因此$T^\vee$成为$\widetilde G$中的分裂环面。

$T^\vee$必为$\widetilde G$的极大环面，因为复化表示环的Krull维数等于$\operatorname{rank}X^\vee$，亦即$T^\vee$的秩。

剩下只需把$\widetilde G$的根数据识别为$G$的对偶根数据。我们已固定正根，等价地固定$T\subseteq B\subseteq G$。已经得到两点：
\begin{itemize}
\item $\widetilde G$的不可约表示为$F(IC_\lambda)$，其标号$\lambda$为$T$的支配余权。
\item Langlands对偶群$G^L$的不可约表示为$V^L_\lambda$，其标号为$T^\vee$的支配权，也就是$T$的支配余权。
\end{itemize}
$V^L_\lambda$中出现的权在Weyl群作用下对称。对$F(IC_\lambda)$，其权集合由半无限轨道与Schubert闭包的交确定，见\eref{笔记}{L21S2O6}，所以与$V^L_\lambda$的权集合相同。
\sourcepage{198}
于是，给$\widetilde G$选择支配室，可以等价地通过为各$F(IC_\lambda)$选择最高权线来完成。选取$F_\lambda(IC_\lambda)$作为最高权线。由最高权与实际出现的权之差确定单根，再由上述权集合的一致性，$\widetilde G$的单根恰为$G$的单余根，故
\[\Phi(\widetilde G,T^\vee)=\Phi(G,T)^\vee.\]
相应的反射及完美配对也识别其余根，因而完整根数据与$G$的对偶根数据一致。因此$\widetilde G\cong G^L$。
\end{proof}
\needspace{5\baselineskip}\section{第23讲参考文献}\label{L23S4}
\begin{references}
\refitem{L23MV}{MV} I. Mirkovi\'c and K. Vilonen, \textit{Geometric Langlands duality and representations of algebraic groups over commutative rings}, Annals of Mathematics, 166 (2007), 95--143.
\refitem{L23JMW}{JMW} D. Juteau, C. Mautner, and G. Williamson, \textit{Parity sheaves}, J. Amer. Math. Soc., 27 (4), 1169--1212, 2014.
\refitem{L23dCM}{dCM} Mark Andrea A. de Cataldo and Luca Migliorini, \textit{The decomposition theorem, perverse sheaves and the topology of algebraic maps}, \S4.4--4.5, available at \url{https://deopurkar.github.io/mhm/dcm.pdf}.
\refitem{L23L}{L} G. Lusztig, \textit{Singularities, character formulas, and a q-analog of weight multiplicities}, \textit{Analysis and topology on singular spaces, II, III} (Luminy, 1981), Ast\'erisque, vol. 101, Soc. Math. France, Paris, 1983, pp. 208--229. MR 737932.
\refitem{L23DM}{DM} P. Deligne and J. Milne, \textit{Tannakian categories}, in \textit{Hodge Cycles and Motives}, Lecture Notes in Math. 900 (1982) 101--228.
\end{references}
